% ADMD Lecture 15. Compile: latexmk -xelatex lecture15.tex
\documentclass[10pt,aspectratio=169,t]{beamer}
\usetheme{upbminimal}

\course{Application Development for Mobile Devices}
\decklabel{Lecture 15}
\title{Release, distribution, and push notifications}
\subtitle{Build modes, signing, the stores, CI/CD, and FCM}
\author{Dinu-Ștefan Rusu}
\institute{FILS, Universitatea Politehnica București}
\date{Autumn 2026}

\begin{document}

\titleframe

\objectivesframe{
  \cell[\faBoxOpen]{Produce a real build}{Debug, profile and release, flavors, icons, and a version number the stores accept.}
  \cell[\faKey]{Sign it correctly}{Android keystores and Play App Signing, iOS certificates and provisioning profiles.}
  \cell[\faRocket]{Deliver it}{Play Console and App Store Connect, internal testing and TestFlight, GitHub Actions and fastlane.}
  \cell[\faBell]{Push to a device}{FCM tokens, the three app states, and what a notification payload may carry.}
}

\divider{Part 1 · Building}{Making an installable build}[Build modes, flavors, icons, and version numbers]

\begin{frame}[eyebrow=Building]{Three build modes, three purposes}
  \vskip 2mm
  \begin{grid}[3]
    \cell[\faBug]{Debug}{JIT compiled, assertions on, hot reload, service extensions, a debug banner. Large and slow on purpose. Never measure performance here.}
    \cell[\faClock]{Profile}{AOT compiled with tracing kept, so DevTools can still see frames and CPU. Runs on a real device only, never on an emulator.}
    \cell[\faRocket]{Release}{AOT compiled, assertions stripped, no debugging hooks, tree shaken icons. This is the only mode a user ever runs.}
  \end{grid}
  \vskip 4mm
  \muted{The mode is a compile-time decision. It changes the Dart compiler, the asserts, and the size of the binary.}
  \begin{takeaway}{A bug that only appears in release is usually an assert or a debug-only code path.}
    Reproduce with \code{flutter run --release} before you blame the store build.
  \end{takeaway}
  \note{Ask which mode they have been running all semester. Then say that everything they measured in DevTools in Lecture 5 was measured in profile mode for exactly this reason.}
\end{frame}

\begin{frame}[fragile,eyebrow=Building]{The artifacts each store wants}
  \begin{columns}[T]
    \begin{column}{0.52\textwidth}
\begin{shell}
# Android: the format Play expects
flutter build appbundle --release
# Sideloading and CI smoke tests only
flutter build apk --split-per-abi
# iOS: an archive, then upload it
flutter build ipa --release
\end{shell}
    \end{column}
    \begin{column}{0.44\textwidth}
      \lead{An app bundle is not an installable app.}{You upload one \code{.aab}, and Play generates the APK for each device: its ABI, screen density and language only.}
      \muted{\code{flutter build ipa} writes an Xcode archive plus an \code{.ipa}. You upload it with Transporter or from Xcode Organizer.}
    \end{column}
  \end{columns}
  \begin{takeaway}{Ship the bundle, keep the debug symbols.}
    Archive the mapping and symbol files from every release build, or your crash reports stay unreadable.
  \end{takeaway}
  \note{Ask who has ever installed an APK by hand. Then point out that Play has not accepted a plain APK for a new app in years. The split APKs are for the lab and for testing on a device you hold.}
\end{frame}

\begin{frame}[fragile,eyebrow=Building]{Flavors: one codebase, several apps}
  \begin{columns}[T]
    \begin{column}{0.54\textwidth}
\begin{shell}
flutter run --flavor dev \
  -t lib/main_dev.dart

flutter build appbundle --flavor prod \
  -t lib/main_prod.dart \
  --dart-define=API_URL=https://api.dev
\end{shell}
      \muted{Each flavor gets its own application id, its own icon and its own name, so development and production can sit side by side on the same phone.}
    \end{column}
    \begin{column}{0.42\textwidth}
      \begin{hint}[Where it is configured]
        Android: product flavors in the Gradle build file. iOS: one scheme and one build configuration per flavor, in Xcode.
      \end{hint}
      \muted{Read the values back with \code{String.fromEnvironment}. A \code{--dart-define} is baked in at compile time, so it is not a secret.}
    \end{column}
  \end{columns}
  \note{The point to make: a flavor changes identity, not code. The staging build must be the same binary logic as production, or staging proves nothing.}
\end{frame}

\begin{frame}[fragile,eyebrow=Building]{Icons, names, and the splash screen}
  \begin{columns}[T]
    \begin{column}{0.5\textwidth}
\begin{yaml}
flutter_launcher_icons:
  image_path: "assets/icon.png"
  android: true
  ios: true
  adaptive_icon_background: "#101010"
  adaptive_icon_foreground: "fg.png"
\end{yaml}
    \end{column}
    \begin{column}{0.46\textwidth}
      \lead{Generate, do not hand-place.}{One square source image, one command, and every density and every iOS size is written for you.}
      \muted{Adaptive icons need a foreground and a background layer, because the launcher masks them into its own shape.}
      \muted{\code{flutter\_native\_splash} writes the native launch screen. Flutter cannot paint before the engine starts.}
    \end{column}
  \end{columns}
  \begin{takeaway}{The app icon may not have a transparent background.}
    iOS rejects the upload for it. Export the source image on a solid color.
  \end{takeaway}
  \note{Warn about the transparent-background icon: iOS does not allow alpha in the app icon and the upload is rejected. Mention that the splash screen is native by necessity, and that a long one is a symptom of slow startup, not a feature.}
\end{frame}

\begin{frame}[fragile,eyebrow=Building]{Version name and build number}
  \begin{columns}[T]
    \begin{column}{0.44\textwidth}
\begin{yaml}
# pubspec.yaml
version: 1.4.0+27
\end{yaml}
\begin{shell}
flutter build appbundle \
  --build-name=1.4.0 --build-number=28
\end{shell}
    \end{column}
    \begin{column}{0.52\textwidth}
      {\usebeamerfont{small}\begin{tabular}{@{}lll@{}}
        \toprule
        \thead{Part} & \thead{Android} & \thead{iOS} \\
        \midrule
        1.4.0 & versionName & CFBundleShortVersionString \\
        27 & versionCode & CFBundleVersion \\
        \bottomrule
      \end{tabular}}
      \vskip 3mm
      \muted{The name is for humans and may repeat. The number is for the store, must be an integer, and must increase with every upload.}
    \end{column}
  \end{columns}
  \begin{takeaway}{A build number is spent the moment you upload it.}
    Both stores refuse a second upload with a number they already have, even if you deleted the first one. Let CI set it from the run number.
  \end{takeaway}
  \note{This is the single most common wasted hour in a first release. Show the pubspec line and say the plus sign separates the two values. Ask what happens when two engineers release from their laptops on the same day.}
\end{frame}

\divider{Part 2 · Signing}{Proving the update is yours}[Keystores, Play App Signing, certificates and profiles]

\begin{frame}[fragile,eyebrow=Signing]{The Android keystore}
  \begin{columns}[T]
    \begin{column}{0.56\textwidth}
\begin{shell}
keytool -genkey -v -keystore upload.jks \
  -keyalg RSA -keysize 2048 -validity 10000 \
  -alias upload
\end{shell}
\begin{yaml}
# android/key.properties  (never committed)
storeFile=/Users/you/keys/upload.jks
storePassword=...
keyAlias=upload
keyPassword=...
\end{yaml}
    \end{column}
    \begin{column}{0.40\textwidth}
      \lead{Every Android app is signed.}{The signature is the identity. Android installs an update only when it is signed by the same key as the installed app.}
      \muted{Gradle reads \code{key.properties} and applies the release signing config. Without it, \code{flutter build} silently falls back to the debug key.}
    \end{column}
  \end{columns}
  \note{Say plainly that a debug-signed bundle is rejected by Play. Point at the gitignore line for key.properties. Ask what a lost key costs: on Android it used to mean a new listing and zero installs.}
\end{frame}

\begin{frame}[eyebrow=Signing]{Play App Signing}
  \vskip 1mm
  \begin{grid}[3]
    \cell[\faUpload]{You hold the upload key}{You sign the bundle with it. Play verifies the signature, then strips it. It only proves the upload came from you.}
    \cell[\faLock]{Google holds the app signing key}{Play re-signs the generated APKs with it. That key is what devices check, and it never leaves Google.}
    \cell[\faIcon{sync-alt}]{The upload key is replaceable}{Lose it and you register a new one through support. Lose a self-managed app signing key and the app can never be updated.}
  \end{grid}
  \vskip 1mm
  \muted{Play App Signing is required for new apps, and it is what makes per-device APK splitting possible.}
  \begin{takeaway}{Back up the keystore and its passwords the day you create them.}
    A password manager entry and an encrypted copy off your laptop. Not the repository.
  \end{takeaway}
  \note{Draw the two keys on the board: upload key at the door, app signing key on the package. Students confuse them constantly.}
\end{frame}

\begin{frame}[eyebrow=Signing]{iOS: four things that must agree}
  {\usebeamerfont{small}\begin{tabular}{@{}lp{56mm}p{30mm}@{}}
    \toprule
    \thead{Piece} & \thead{What it is} & \thead{Where it lives} \\
    \midrule
    Certificate & Apple's signature on your public key & Keychain, exported as \code{.p12} \\
    App ID & The bundle identifier and its capabilities & Developer portal \\
    Devices & UDIDs allowed to run a test build & Developer portal \\
    Profile & Certificate, App ID and devices, in one file & \code{.mobileprovision} \\
    \bottomrule
  \end{tabular}}
  \vskip 2mm
  \muted{Xcode manages all four for one developer. A team shares them deliberately: \code{fastlane match} keeps them in a private encrypted repository.}
  \begin{takeaway}{Certificates expire.}
    A build that worked last year fails today for no other reason.
  \end{takeaway}
  \note{This table is the answer to almost every iOS signing error message. Tell students to read the error and ask which of the four rows disagrees with the others.}
\end{frame}

\begin{frame}[eyebrow=Signing]{What never goes into the repository}
  \vskip 2mm
  \begin{columns}[T]
    \begin{column}{0.48\textwidth}
      \begin{warning}[Keep these out of Git]
        Keystores and \code{key.properties}. Signing certificates and \code{.p12} exports. Provisioning profiles. App Store Connect API keys. Play service account JSON. Any \code{.env} with a real secret.
      \end{warning}
    \end{column}
    \begin{column}{0.48\textwidth}
      \lead{A secret in Git is public.}{Deleting the file does not help: the object stays in history, and on a public repository it is scraped within minutes.}
      \muted{If it happens: rotate the credential first, then clean the history. In that order. Rotation is what actually stops the damage.}
      \muted{Store the real values in the CI provider's encrypted secrets and in a password manager for humans.}
    \end{column}
  \end{columns}
  \begin{takeaway}{Write the ignore rules before you create the keys.}
    That is the only moment when the repository is still guaranteed clean.
  \end{takeaway}
  \note{Mention that a bot scanning new public commits is faster than any of us. Ask the room to check their own lab repository for a stray key before the lab test.}
\end{frame}

\divider{Part 3 · Stores}{Getting past review}[Play Console, App Store Connect, test tracks, and privacy]

\begin{frame}[eyebrow=Stores]{Google Play Console: the first release}
  \vskip 1mm
  \begin{enumerate}
    \item Create the app and fix the package name. \muted{It is permanent.}
    \item Complete the store listing. \muted{Descriptions, icon, feature graphic, screenshots.}
    \item Answer the questionnaires. \muted{Content rating, target audience, ads, Data safety.}
    \item Upload the signed bundle to a track, with release notes.
    \item Roll out in stages, and watch the crash rate as it climbs.
  \end{enumerate}
  \vskip 3mm
  \muted{Review takes days, not minutes, and a new developer account waits longest.}
  \note{The realistic message: the binary is the easy half. Budget an afternoon for the forms and screenshots, and do it before the deadline, not on it.}
\end{frame}

\begin{frame}[eyebrow=Stores]{App Store Connect: the first release}
  \vskip 1mm
  \begin{enumerate}
    \item Register the App ID and create the app record. \muted{Bundle identifier and SKU, both permanent.}
    \item Upload a build. \muted{From Xcode Organizer or Transporter. It appears after processing.}
    \item Fill in the metadata. \muted{Screenshots per device size, description, support and privacy URLs.}
    \item Answer App Privacy and export compliance, attach the build, submit.
    \item Release manually or in phases once it is approved.
  \end{enumerate}
  \vskip 3mm
  \muted{Give the reviewer a demo account and notes for anything behind a login. A reviewer who cannot get in rejects the app.}
  \note{Say that Apple's review is a human reading your notes and tapping through the app. Everything you would tell a new tester belongs in the review notes field.}
\end{frame}

\begin{frame}[eyebrow=Stores]{Test tracks before anyone real sees it}
  \vskip 1mm
  {\usebeamerfont{small}\begin{tabular}{@{}llp{40mm}p{40mm}@{}}
    \toprule
    \thead{Track} & \thead{Store} & \thead{Who gets it} & \thead{Review} \\
    \midrule
    Internal testing & Play & a small list of accounts you name & minimal, available quickly \\
    Closed testing & Play & testers on a list or in a group & yes \\
    Open testing & Play & anyone with the opt-in link & yes \\
    TestFlight internal & App Store & members of your team & none \\
    TestFlight external & App Store & invited testers or a public link & Beta App Review \\
    \bottomrule
  \end{tabular}}
  \vskip 3mm
  \begin{takeaway}{Every release goes through a test track first.}
    A build that installs and launches on ten real devices has already caught the failures an emulator hides.
  \end{takeaway}
  \note{Point out that TestFlight builds expire after a while, so a beta is not a distribution channel. Internal testing on Play is the fastest loop you have.}
\end{frame}

\begin{frame}[eyebrow=Stores]{Why apps get rejected}
  \vskip 1mm
  \begin{grid}[3]
    \cell[\faBug]{It crashed for the reviewer}{The most common reason. Test the exact build you submitted.}
    \cell[\faLock]{The reviewer could not sign in}{No demo account, a dead link, or a phone code they cannot receive.}
    \cell[\faFile]{Placeholder content}{Lorem ipsum, dead links, or a support URL that does not resolve.}
    \cell[\faIcon{user-shield}]{Permissions without a reason}{A permission you request but never use.}
    \cell[\faSearch]{Misleading metadata}{Screenshots of features that do not exist.}
    \cell[\faMobile]{Too little to be an app}{A thin wrapper around a website.}
  \end{grid}
  \note{Tell them a rejection is a message, not a verdict: you fix it and resubmit the same day. The expensive mistake is submitting untested on the last day.}
\end{frame}

\begin{frame}[eyebrow=Stores]{Privacy declarations and permissions}
  \begin{columns}[T]
    \begin{column}{0.48\textwidth}
      \lead{You declare data before you collect it.}{Play calls it the Data safety form, Apple calls them privacy labels. Both ask what you collect, why, and whether you share it.}
      \muted{The declaration must match the app, including whatever your analytics and crash SDKs send. An SDK that collects an advertising identifier is your collection, not theirs.}
    \end{column}
    \begin{column}{0.48\textwidth}
      \begin{hint}[On the device]
        iOS needs a purpose string in \code{Info.plist} for every sensitive API, shown verbatim in the prompt. Android declares permissions in the manifest and asks at runtime.
      \end{hint}
      \muted{Ask at the moment the permission is needed, with the reason on screen first. A cold prompt at launch is denied, and a denial is hard to reverse.}
    \end{column}
  \end{columns}
  \note{Connect this back to Lecture 11 on permissions. The new part here is that the store now asks you to write down, in public, what Lecture 9 and Lecture 12 actually send off the device.}
\end{frame}

\divider{Part 4 · Automation}{Nobody releases from a laptop}[GitHub Actions, fastlane, Codemagic, and code push]

\begin{frame}[eyebrow=Automation]{What a release pipeline does}
  \vskip 2mm
  \begin{grid}[3]
    \cell[\faCheckCircle]{On every pull request}{Format check, \code{flutter analyze}, unit and widget tests. Merging is blocked while any of them is red.}
    \cell[\faBoxOpen]{On every merge to main}{Build the release artifacts for both platforms so a broken build is found in minutes, not at release time.}
    \cell[\faRocket]{On a version tag}{Sign, set the build number, upload to internal testing and TestFlight, and post the link to the team channel.}
  \end{grid}
  \vskip 4mm
  \muted{The pipeline is also the honest answer to "it works on my machine". It builds from a clean checkout every time, with no local caches and no local keys.}
  \begin{takeaway}{Automate the boring half first.}
    Tests and analysis on pull requests pay for themselves in a week. Automated store upload can wait until you release often enough to care.
  \end{takeaway}
  \note{Callback to Lecture 14: the tests written there are what this pipeline runs. A test suite nobody runs automatically decays within a month.}
\end{frame}

\begin{frame}[fragile,eyebrow=Automation]{A GitHub Actions workflow}
\begin{yaml}
name: ci
on: [pull_request]
jobs:
  test:
    runs-on: ubuntu-latest
    steps:
      - uses: actions/checkout@v4
      - uses: subosito/flutter-action@v2
      - run: flutter pub get
      - run: dart format --set-exit-if-changed .
      - run: flutter analyze
      - run: flutter test
\end{yaml}
  \muted{Format, analyze and test are what you run locally: a red pipeline is never a surprise.}
  \note{Show that the whole file is about ten meaningful lines. Say that an iOS build needs a macOS runner, which is slower, so that job belongs on a tag and not on every pull request.}
\end{frame}

\begin{frame}[fragile,eyebrow=Automation]{Secrets in the pipeline}
  \begin{columns}[T]
    \begin{column}{0.54\textwidth}
\begin{yaml}
- name: Restore the keystore
  env:
    KEYSTORE: ${{ secrets.KEYSTORE_B64 }}
  run: |
    echo "$KEYSTORE" | base64 -d \
      > android/upload.jks
\end{yaml}
    \end{column}
    \begin{column}{0.42\textwidth}
      \lead{Binary secrets travel as base64.}{Store them in the provider's encrypted secrets, decode them into the workspace, and let the job delete the workspace afterwards.}
      \muted{Never echo a secret, never print it in a failing step, and never expose secrets to a workflow triggered by a fork's pull request.}
    \end{column}
  \end{columns}
  \begin{takeaway}{A secret printed once in a build log is a burned secret.}
    Rotate it. Deleting the log does not help.
  \end{takeaway}
  \note{Point at the pipe into base64 minus d and note there is no echo of the value itself. A secret printed once in a public build log is a rotated secret.}
\end{frame}

\begin{frame}[eyebrow=Automation]{fastlane and Codemagic}
  {\usebeamerfont{small}\begin{tabular}{@{}lp{48mm}p{48mm}@{}}
    \toprule
    \thead{} & \thead{fastlane} & \thead{Codemagic} \\
    \midrule
    What it is & Ruby lanes you run anywhere & Hosted CI built around Flutter \\
    Signing & \code{match} shares certificates through a private repo & Managed in the UI, on a machine with Xcode \\
    Upload & \code{supply} for Play, \code{pilot} and \code{deliver} for Apple & Publishing to both stores is a built-in step \\
    Fits & teams that want the pipeline in their repo & teams that do not want to run macOS runners \\
    \bottomrule
  \end{tabular}}
  \vskip 2mm
  \muted{They are not exclusive: a common setup is GitHub Actions for tests and fastlane for everything that talks to a store.}
  \note{Say that both save the same thing: the twenty manual clicks between a green build and a build in testers' hands. Pick one, write it once, forget it.}
\end{frame}

\begin{frame}[eyebrow=Automation]{Code push, in one slide}
  \begin{columns}[T]
    \begin{column}{0.5\textwidth}
      \lead{Shorebird patches Dart code on installed apps.}{You publish a patch, the app downloads it, and the fix is live on the next launch without a store review.}
      \muted{It covers Dart code and assets. Anything native, a new plugin, a new permission, or a version bump still needs a real store release.}
    \end{column}
    \begin{column}{0.46\textwidth}
      \begin{warning}[Use it narrowly]
        A patch must stay inside what the stores allow: a fix to the app you shipped, not a new product. Treat it as a hotfix channel, and keep every patch in Git.
      \end{warning}
      \muted{It also splits your users across two code paths. Know which patch a crash report came from before you debug it.}
    \end{column}
  \end{columns}
  \note{Frame the trade-off honestly: hours instead of days for a critical fix, at the price of another moving part. Most teams never need it.}
\end{frame}

\divider{Part 5 · Push}{Waking an app that is closed}[FCM, APNs, tokens, and the three app states]

\begin{frame}[eyebrow=Push]{How a push actually arrives}
  \vskip 2mm
  \begin{grid}[3]
    \cell[\faServer]{Your backend decides}{It calls the FCM HTTP v1 API with a service account credential. The app never sends a push, because the app cannot hold that key.}
    \cell[\faCloud]{FCM routes it}{Firebase Cloud Messaging picks the transport: Google Play services on Android, APNs on iOS. One API, two networks.}
    \cell[\faMobile]{The device delivers it}{The system wakes your app or draws the notification itself, depending on the payload and on what the app is doing.}
  \end{grid}
  \vskip 4mm
  \muted{Delivery is best effort. Messages can be delayed by battery saving, collapsed with a newer one, or dropped while the device is offline.}
  \begin{takeaway}{A push is a hint, not a transport.}
    Send an identifier and let the app fetch the real data over HTTP. Never treat a push as your only copy of a state change.
  \end{takeaway}
  \note{Ask why the app cannot send its own pushes. The answer is the same as the API key trap in Lecture 12: anything in the binary belongs to the attacker too.}
\end{frame}

\begin{frame}[fragile,eyebrow=Push]{Setup and the registration token}
  \begin{columns}[T]
    \begin{column}{0.54\textwidth}
\begin{dart}
final fm = FirebaseMessaging.instance;
await fm.requestPermission();

final token = await fm.getToken();
await api.registerDevice(token);

fm.onTokenRefresh.listen(api.register);
\end{dart}
    \end{column}
    \begin{column}{0.42\textwidth}
      \lead{The token is the address.}{One token per app install. It changes on reinstall, on a data clear, and sometimes on its own.}
      \muted{Store it against the user and the device, delete it on sign-out, and drop it when the send API reports it unregistered.}
      \muted{iOS also needs the APNs auth key in the Firebase project, plus push and background modes enabled in Xcode.}
    \end{column}
  \end{columns}
  \note{Stress the refresh stream: the classic bug is registering the token once at first launch and never again, so a fraction of users silently stops receiving anything.}
\end{frame}

\begin{frame}[fragile,eyebrow=Push]{Notification messages and data messages}
  \begin{columns}[T]
    \begin{column}{0.46\textwidth}
\begin{yaml}
message:
  token: "..."
  notification:
    title: "New comment"
    body: "Ana replied to your task"
  data:
    taskId: "8412"
\end{yaml}
    \end{column}
    \begin{column}{0.5\textwidth}
      \muted{A \code{notification} block is drawn by the operating system when your app is not in the foreground. Your code does not run until the user taps it.}
      \muted{A \code{data} block is never drawn. It is handed to your app, so you decide what to show. On iOS it needs the background flag and is delivered at the system's convenience.}
      \lead{Sending both is the usual choice.}{The system shows something reliable, and the data payload tells you which screen to open on tap.}
    \end{column}
  \end{columns}
  \note{The pitfall to name: data values are strings only. Anything structured has to be encoded, and the whole payload is capped at a few kilobytes.}
\end{frame}

\begin{frame}[fragile,eyebrow=Push]{Handling the three app states}
  \begin{columns}[T]
    \begin{column}{0.56\textwidth}
\begin{dart}
@pragma('vm:entry-point')
Future<void> bg(RemoteMessage m) async {
  await Firebase.initializeApp();
}
FirebaseMessaging.onBackgroundMessage(bg);

fm.onMessage.listen(showLocal);
fm.onMessageOpenedApp.listen(openScreen);
final m = await fm.getInitialMessage();
\end{dart}
    \end{column}
    \begin{column}{0.40\textwidth}
      \muted{Foreground: nothing is drawn for you. Show a local notification, or update the screen.}
      \muted{Background: the system draws it, and \code{onMessageOpenedApp} fires on tap.}
      \muted{Terminated: the app starts cold, so \code{getInitialMessage} is the only place the tapped message appears.}
      \begin{hint}[Background handler]
        Top level or static, annotated, and it initializes Firebase itself.
      \end{hint}
    \end{column}
  \end{columns}
  \note{Make them test all three by hand: app open, app backgrounded, app force stopped. Almost every push bug reported in the labs is one of the three cases never being tried.}
\end{frame}

\begin{frame}[eyebrow=Push]{Asking for permission, and what to send}
  \begin{columns}[T]
    \begin{column}{0.48\textwidth}
      \lead{You get one prompt.}{iOS has always required permission. Recent Android versions require it too. A denial is final until the user walks into system settings.}
      \muted{Ask after the user has done something that a notification would follow: created a task with a due date, joined a conversation. Explain the value on your own screen first.}
      \muted{Group related messages, collapse duplicates, and respect quiet hours. An app that pushes daily noise gets its notifications switched off, then deleted.}
    \end{column}
    \begin{column}{0.48\textwidth}
      \begin{warning}[Not in a payload]
        Passwords, tokens, one-time codes, health or financial details, anything private enough to matter on a lock screen. The payload passes through Google and Apple, and is stored on the device.
      \end{warning}
      \muted{Send an identifier and a neutral title. Fetch the sensitive content after the user opens the app and is authenticated.}
    \end{column}
  \end{columns}
  \note{Ask the room how many have disabled notifications for an app this month, and why. That answer is the design brief for this slide.}
\end{frame}

\begin{frame}[eyebrow=Checklist]{Before you ship}
  \vskip 2mm
  \begin{columns}[T]
    \begin{column}{0.48\textwidth}
      \begin{checklist}
        \item Release build tested on a physical device, not only an emulator
        \item Build number increased, release notes written
        \item Signing keys backed up outside the repository
        \item Crash reporting and analytics live in the release build
      \end{checklist}
    \end{column}
    \begin{column}{0.48\textwidth}
      \begin{checklist}
        \item Privacy declarations match what the app actually sends
        \item Every permission has a purpose string and a reason on screen
        \item A demo account and review notes are attached
        \item A staged rollout, and someone watching the crash rate
      \end{checklist}
    \end{column}
  \end{columns}
  \vskip 3mm
  \muted{Run this list from a written checklist, not from memory. The one you skip is the one that fails.}
  \note{Suggest they copy this into the pull request template of their own project. A release checklist is the cheapest quality tool that exists.}
\end{frame}

\begin{frame}[eyebrow=Course recap]{Fourteen weeks in one table}
  \vskip 1mm
  \begin{columns}[T]
    \begin{column}{0.48\textwidth}
      {\usebeamerfont{small}\begin{tabular}{@{}ll@{}}
        \toprule
        \thead{Wk} & \thead{Topic} \\
        \midrule
        1 & Orientation, Git, build strategies \\
        2 & Dart, null safety, toolchain \\
        3 & Widgets and layout \\
        4 & Async Dart, coding agents \\
        5 & Debugging, state part 1 \\
        6 & BLoC and Riverpod \\
        7 & HTTP, JSON, code generation \\
        \bottomrule
      \end{tabular}}
    \end{column}
    \begin{column}{0.48\textwidth}
      {\usebeamerfont{small}\begin{tabular}{@{}ll@{}}
        \toprule
        \thead{Wk} & \thead{Topic} \\
        \midrule
        8 & Firebase, flags, GraphQL \\
        9 & Observability, local storage \\
        10 & Auth, OAuth, App Check \\
        11 & Routing, permissions, sockets \\
        12 & AI in the app \\
        13 & UI polish and testing \\
        14 & Release, distribution, push \\
        \bottomrule
      \end{tabular}}
    \end{column}
  \end{columns}
  \note{Walk the table quickly and name the lab that used each week. The point is that the todo app in their repository touched every row.}
\end{frame}

\begin{frame}[eyebrow=Lab test]{The lab test and your grade}
  \begin{columns}[T]
    \begin{column}{0.5\textwidth}
      \lead{Lab 7 is the practical test.}{Individual, 90 minutes, your own laptop, internet allowed. Agents are allowed, and you must be able to explain every line you submit.}
      \muted{A one-page spec of a small app with four requirements: a screen built from the layout widgets, state in a Cubit, one network call with a model class, one navigation with \code{go\_router}.}
      \muted{One point per requirement, partial credit for a requirement that is present but incomplete. Submit on branch \code{lab-7} before the end.}
    \end{column}
    \begin{column}{0.46\textwidth}
      {\usebeamerfont{small}\begin{tabular}{@{}lc@{}}
        \toprule
        \thead{Component} & \thead{Points} \\
        \midrule
        Lab assignments, from the repo & 5 \\
        Practical lab test, Lab 7 & 4 \\
        Participation & 1 \\
        \midrule
        Pass & 5 of 10 \\
        \bottomrule
      \end{tabular}}
      \vskip 3mm
      \muted{Every requirement maps to a lab you already did. Nothing in the test is new material.}
    \end{column}
  \end{columns}
  \note{Read the four requirements out loud and name the lab that taught each one. Then say the honest preparation advice: rebuild Lab 5 from an empty project in an hour, and the test is straightforward.}
\end{frame}

\begin{frame}[eyebrow=Recap]{Three things to remember}
  \vskip 2mm
  \begin{enumerate}
    \item Release is a process, not a build command. \muted{Signing, versioning, forms and review take longer than the compile.}
    \item Guard the keys, automate the rest. \muted{A lost signing key is unrecoverable. Everything else can be scripted.}
    \item A push is a hint. \muted{Send an identifier, fetch the content, and never put a secret in a payload.}
  \end{enumerate}
  \vskip 5mm
  \kv{Further reading}{\link{https://docs.flutter.dev/deployment/android}{docs.flutter.dev/deployment/android} · \link{https://docs.flutter.dev/deployment/ios}{deployment/ios} · \link{https://firebase.google.com/docs/cloud-messaging}{Firebase Cloud Messaging} · \link{https://docs.fastlane.tools}{docs.fastlane.tools}}
  \note{Close by pointing at the two Flutter deployment pages: they are short, current, and they are the checklist for their first real release.}
\end{frame}

\closingframe{Questions?}{dinu\_stefan.rusu@upb.ro · Microsoft Teams: course channel}

\end{document}
